\section{Steady-State Error and Stability of Digital Control Systems}\label{sec:analysis}
\lecdate{11, 18 and 25 November 2019}
% Figures of Chapter 4 (steady-state error and stability), redrawn from the board sketches.

% ---- p.43: s-plane and z-plane stability regions ----------------------------------
\newcommand{\figSZplanes}{%
\begin{tikzpicture}
  % s plane
  \fill[pattern=north east lines,pattern color=gray!55] (-2.1,-1.7) rectangle (0,1.8);
  \draw[->] (-2.1,0) -- (2.0,0) node[right,lab]{$\Real[s]$};
  \draw[->] (0,-1.7) -- (0,1.95) node[above,lab]{$\Imag[s]$};
  % z plane
  \begin{scope}[xshift=7.2cm]
    \fill[pattern=north east lines,pattern color=gray!55] (0,0) circle (1.4);
    \draw[thick] (0,0) circle (1.4);
    \draw[->] (-2.1,0) -- (2.2,0) node[right,lab]{$\Real[z]$};
    \draw[->] (0,-1.7) -- (0,1.95) node[above,lab]{$\Imag[z]$};
    \node[lab,below left] at (-1.4,0) {$-1$};
    \node[lab,below right] at (1.4,0) {$1$};
  \end{scope}
\end{tikzpicture}}

% ---- p.47: bounded output sequence ------------------------------------------------
\newcommand{\figBiboSeq}{%
\begin{tikzpicture}
  \begin{axis}[width=0.62\linewidth,height=4.2cm,axis lines=middle,xlabel={$k$},ylabel={$y(k)$},
      xmin=-0.5,xmax=15.5,ymin=-1.5,ymax=1.5,xtick=\empty,ytick={-1,1},yticklabels={$-b_y$,$b_y$},
      every axis x label/.style={at={(ticklabel* cs:1.0)},anchor=west},
      every axis y label/.style={at={(ticklabel* cs:1.0)},anchor=south}]
    \addplot[ycomb,itublue,thick,mark=*,mark size=1.5pt] coordinates
      {(1,0.15)(2,0.45)(3,-0.35)(4,0.1)(5,0.3)(6,-0.6)(7,0.7)(8,-0.2)(9,0.9)(10,-0.5)(11,-0.85)(12,0.4)(13,0.55)(14,-0.3)};
    \draw[dashed,gray] (axis cs:-0.5,1) -- (axis cs:15.5,1);
    \draw[dashed,gray] (axis cs:-0.5,-1) -- (axis cs:15.5,-1);
  \end{axis}
\end{tikzpicture}}

% ---- p.48: unit circle with poles and the unbounded ramp ---------------------------
\newcommand{\figPoleCircle}{%
\begin{tikzpicture}
  \draw[thick] (0,0) circle (1.4);
  \draw[->] (-1.9,0) -- (2.1,0) node[right,lab]{$\Real[z]$};
  \draw[->] (0,-1.8) -- (0,1.9) node[above,lab]{$\Imag[z]$};
  \node[lab,below left] at (-1.4,0) {$-1$};
  \node[lab,below right] at (1.4,0) {$1$};
  \node[pole] at (0.5,0) {};
  \node[pole,draw=boardgreen] at (1.4,0) {};
  \node[pole] at (-0.4,0.6) {}; \node[pole] at (-0.4,-0.6) {};
\end{tikzpicture}}

\newcommand{\figRampSeq}{%
\begin{tikzpicture}
  \begin{axis}[width=0.42\linewidth,height=3.8cm,axis lines=middle,xlabel={$k$},ylabel={$y(k)$},
      xmin=-0.3,xmax=6.5,ymin=0,ymax=6,xtick=\empty,ytick=\empty,
      every axis x label/.style={at={(ticklabel* cs:1.0)},anchor=west},
      every axis y label/.style={at={(ticklabel* cs:1.0)},anchor=south}]
    \addplot[ycomb,itublue,thick,mark=*,mark size=1.5pt,samples at={0,1,...,5}]{x};
  \end{axis}
\end{tikzpicture}}

% ---- p.50: BIBO regions -------------------------------------------------------------
\newcommand{\figBiboRegions}{%
\begin{tikzpicture}
  \draw[dotted,thick] (0,0) circle (1.5);
  \draw[->] (-2.2,0) -- (2.4,0) node[right,lab]{$\Real[z]$};
  \draw[->] (0,-1.9) -- (0,2.0) node[above,lab]{$\Imag[z]$};
  \node[lab,boardgreen] at (0,0.8) {BIBO stable};
  \node[lab,boardgreen,align=left] (u) at (2.5,1.9) {BIBO\\ unstable};
  \draw[->,boardgreen] (u.south west) -- (1.3,1.0);
\end{tikzpicture}}

% ---- p.53: geometric check of the bilinear transformation --------------------------
\newcommand{\figBilinear}{%
\begin{tikzpicture}[scale=1.0]
  \def\R{2.2}
  % z-plane
  \draw[thick,gray] (0,0) circle (\R);
  \draw[->] (-2.8,0) -- (2.9,0) node[right,lab]{$\Real[z]$};
  \draw[->] (0,-2.5) -- (0,2.7) node[above,lab]{$\Imag[z]$};
  \coordinate (m1) at (-\R,0); \coordinate (p1) at (\R,0);
  \coordinate (z1) at (0.88,-1.43); \coordinate (z2) at (1.32,1.76); \coordinate (z3) at (-1.98,1.98);
  \foreach \p/\c in {z1/itublue,z2/boardgreen,z3/gray}{
    \draw[\c,thick] (m1) -- (\p); \draw[\c,thick] (p1) -- (\p);
    \fill[\c] (\p) circle (1.5pt);
  }
  \node[lab,itublue,below right] at (z1) {$z_1$};
  \node[lab,boardgreen,above right] at (z2) {$z_2$};
  \node[lab,gray,above left] at (z3) {$z_3$};
  \node[lab,below left] at (m1) {$-1$}; \node[lab,below right] at (p1) {$1$};
  % right angle at z2
  \draw[boardgreen] ($(z2)+(-0.27,-0.13)$) -- ++(0.134,-0.268) -- ++(0.268,0.134);
  % w-plane
  \begin{scope}[xshift=7.6cm]
    \draw[->] (-2.2,0) -- (2.4,0) node[right,lab]{$\Real[w]$};
    \draw[->] (0,-1.7) -- (0,3.4) node[above,lab]{$\Imag[w]$};
    \fill[itublue] (-0.23,-0.71) circle (1.5pt) node[below left,lab]{$w_1$};
    \fill[boardgreen] (0,0.65) circle (1.5pt) node[left,lab]{$w_2$};
    \fill[gray] (0.98,2.85) circle (1.5pt) node[right,lab]{$w_3$};
    \node[lab,align=left] at (1.5,-1.2) {$w=\dfrac{z-1}{z+1}$};
  \end{scope}
\end{tikzpicture}}

% ---- Jury table (p.57) --------------------------------------------------------------
\newcommand{\tabJury}{%
\renewcommand{\arraystretch}{1.35}
\small
\begin{tabular}{c|cccccccc}
\textbf{Row} & $z^{0}$ & $z^{1}$ & $z^{2}$ & $\cdots$ & $z^{n-k}$ & $\cdots$ & $z^{n-1}$ & $z^{n}$\\\hline
1 & $a_0$ & $a_1$ & $a_2$ & $\cdots$ & $a_{n-k}$ & $\cdots$ & $a_{n-1}$ & $a_n$\\
2 & $a_n$ & $a_{n-1}$ & $a_{n-2}$ & $\cdots$ & $a_k$ & $\cdots$ & $a_1$ & $a_0$\\
3 & $b_0$ & $b_1$ & $b_2$ & $\cdots$ & $b_{n-k}$ & $\cdots$ & $b_{n-1}$ & \\
4 & $b_{n-1}$ & $b_{n-2}$ & $b_{n-3}$ & $\cdots$ & $b_{k-1}$ & $\cdots$ & $b_0$ & \\
5 & $c_0$ & $c_1$ & $c_2$ & $\cdots$ & $c_{n-2}$ & & & \\
6 & $c_{n-2}$ & $c_{n-3}$ & $c_{n-4}$ & $\cdots$ & $c_0$ & & & \\
$\vdots$ & $\vdots$ & $\vdots$ & $\vdots$ & & & & & \\
$2n{-}5$ & $s_0$ & $s_1$ & $s_2$ & $s_3$ & & & & \\
$2n{-}4$ & $s_3$ & $s_2$ & $s_1$ & $s_0$ & & & & \\
$2n{-}3$ & $r_0$ & $r_1$ & $r_2$ & & & & &
\end{tabular}}


%==============================================================================
\subsection{Steady-state error and error constants}\scanref{pp.\,42--46}
We consider the unity-feedback block diagram of Fig.~\ref{fig:ufb}.

\begin{figure}[ht]
\centering
\resizebox{0.7\linewidth}{!}{\figEquivLoop{$R(z)$}{$Y(z)$}}
\caption{Unity-feedback digital control system.}\label{fig:ufb}
\end{figure}

From the diagram, $Y(z)=E(z)\,C(z)\,\GZAS(z)$ and $E(z)=R(z)-Y(z)$, hence
\begin{equation}
  E(z)=\frac{R(z)}{1+\GZAS(z)C(z)}=\frac{R(z)}{1+L(z)},\qquad L(z)=\GZAS(z)\,C(z)\quad\text{(loop gain).}
\end{equation}
By the final value theorem, the steady-state error is
\begin{equation}\label{eq:ess-star}
  e(\infty)=\lim_{z\to1}(1-z^{-1})E(z)=\frac{z-1}{z}\cdot\frac{R(z)}{1+L(z)}\Big|_{z\to1}
  =\frac{(z-1)\,R(z)}{z\,\bigl(1+L(z)\bigr)}\Big|_{z=1}.\tag{$*$}
\end{equation}
The limit exists if all $(z-1)$ terms in the denominator cancel. This depends on the reference input
$R(z)$ and on the loop gain $L(z)$.

\paragraph{System type.} We can rewrite $L(z)$ in the form
\begin{equation}
  L(z)=\frac{N(z)}{(z-1)^{n}\,D(z)},
\end{equation}
where $N(z)$ and $D(z)$ are polynomials without unity roots $(z-1)$.

\begin{definition}[Type number]
The type number $n$ is the number of unity poles $(z-1)$ in the $z$ transfer function.
\end{definition}
For example, $G(z)=\dfrac{z+2}{z(z-1)}$ is a type~1 transfer function.

The factor $(z-1)$, i.e.\ $z=1$, plays the same role as a pole at the origin of an $s$-domain transfer
function in determining the steady-state error of the system. The reason is the mapping $z=\ee^{sT}$:
$\dfrac1s\leftrightarrow\dfrac{1}{z-1}$, and $s=0$ corresponds to $z=\ee^{0}=1$.

\begin{figure}[ht]
\centering
\resizebox{0.78\linewidth}{!}{\figSZplanes}
\caption{Stable regions (hatched): the left half of the $s$ plane and the interior of the unit circle in the
$z$ plane.}
\end{figure}

\begin{boardnote}
When $L(z)$ is substituted in $(*)$ the board writes the denominator as $z\bigl(N(z)+(z+1)^{n}D(z)\bigr)$.
The factor is $(z-1)^{n}$, as the derivation shows, and it is typeset that way below.
\end{boardnote}
Substituting $L(z)$ in the expression $(*)$ gives
\begin{equation}
  e_{ss}=e(\infty)=\frac{(z-1)^{n+1}\,D(z)\,R(z)}{z\bigl(N(z)+(z-1)^{n}D(z)\bigr)}\Big|_{z=1}.
\end{equation}

\subsubsection*{Sampled step input}
$\dfrac1s\to R(z)=\dfrac{z}{z-1}$, so
\[
  e_{ss}=\frac{(z-1)\cdot z}{z(z-1)(1+L(z))}\Big|_{z=1}=\frac{1}{1+L(1)}=\frac{1}{1+K_p}.
\]
$K_p=L(1)$ is the \emph{position error constant}: $K_p=\dfrac{N(z)}{(z-1)^{n}D(z)}\Big|_{z=1}=L(1)$ and
\[
  e(\infty)=\begin{cases}\dfrac{1}{1+L(1)},&n=0,\\[2mm] 0,&n\ge1.\end{cases}
\]

\subsubsection*{Sampled ramp input}
$\dfrac{1}{s^{2}}\to R(z)=\dfrac{Tz^{-1}}{(1-z^{-1})^{2}}=\dfrac{Tz}{(z-1)^{2}}$, so
\[
  e_{ss}=\frac{(z-1)R(z)}{z(1+L(z))}\Big|_{z=1}=\frac{T}{(z-1)\bigl(1+L(z)\bigr)}\Big|_{z=1}=\frac{1}{K_v},
\]
with the \emph{velocity error constant}
\[
  K_v=\frac1T\,(z-1)L(z)\Big|_{z=1},\qquad
  e(\infty)=\begin{cases}\infty,&n=0,\\[1mm]\dfrac{T}{(z-1)L(z)\big|_{z=1}},&n=1,\\[3mm]0,&n\ge2.\end{cases}
\]

\subsubsection*{Sampled parabolic input}
The \emph{acceleration constant} is $K_a=\dfrac{1}{T^{2}}(z-1)^{2}L(z)\Big|_{z=1}$ and
\[
  e(\infty)=\begin{cases}\infty,&n\le1,\\[1mm]\dfrac{T^{2}}{(z-1)^{2}L(z)\big|_{z=1}},&n=2,\\[3mm]0,&n\ge3.\end{cases}
\]
\begin{remark}[added]
The board does not write the parabolic input. The formulas above correspond to $r(t)=t^{2}/2$, for which
$R(z)=\dfrac{T^{2}z(z+1)}{2(z-1)^{3}}$; then $e_{ss}=\dfrac{T^{2}(z+1)}{2(z-1)^{2}(1+L(z))}\Big|_{z=1}=\dfrac1{K_a}$.
All of this assumes that the closed loop is stable; otherwise the final value theorem does not apply.
\end{remark}

\begin{table}[ht]
\centering
\caption{Steady-state error $e(\infty)$ of a unity-feedback digital system (summary of the formulas above).}\label{tab:ess}
\renewcommand{\arraystretch}{1.5}
\begin{tabular}{c|c|c|c}
Type $n$ & Step & Ramp & Parabola\\\hline
0 & $\dfrac{1}{1+K_p}$ & $\infty$ & $\infty$\\
1 & $0$ & $\dfrac{1}{K_v}$ & $\infty$\\
2 & $0$ & $0$ & $\dfrac{1}{K_a}$\\
$\ge3$ & $0$ & $0$ & $0$
\end{tabular}
\end{table}

\begin{example}[Type-1 loop, step and ramp inputs]
Find the steady-state error $e(\infty)$ for the digital control system with unity feedback and transfer
functions
\[
  \GZAS(z)=\frac{K(z+a)}{(z-1)(z-b)},\qquad C(z)=\frac{K_c(z-b)}{z-c},\qquad 0<a,b,c<1,
\]
(a) for a sampled unit step and (b) for a sampled unit ramp.

\noindent The loop gain is
\[
  L(z)=C(z)\GZAS(z)=\frac{K K_c\,(z+a)}{(z-1)(z-c)},\qquad L(1)=\infty,
\]
so the system type is $n=1$.
\begin{itemize}
  \item For a sampled unit step, $e(\infty)=\dfrac{1}{1+L(1)}=0$.
  \item For a sampled unit ramp,
  \[
    e(\infty)=\frac{T}{(z-1)L(z)\big|_{z=1}}=\frac{T}{\dfrac{K K_c(1+a)}{1-c}}=\frac{T(1-c)}{K K_c(1+a)}.
  \]
\end{itemize}
A larger $K_c$ gives a smaller $e_{ss}$ ($K_c\uparrow\Rightarrow e_{ss}\downarrow$). The controller poles and zeros
affect the steady-state error.
\end{example}

\begin{example}[First-order plant with a proportional controller]\label{ex:ess-first}
Find the steady-state error $e(\infty)$ for the analog system $G(s)=\dfrac{K}{s+a}$ with a proportional
digital controller $K_c$ and a sampled unit-step input. The ZOH, $G(s)$ and the sampler form the block
$\GZAS(z)=(1-z^{-1})\ZT{G(s)/s}$ of the loop of Fig.~\ref{fig:ufb}, with $C(z)=K_c$ and
$L(z)=K_c\GZAS(z)$.
\[
  G(s)=\frac{K}{s+a}\ \longrightarrow\ \GZAS(z)=\frac{z-1}{z}\,\ZT{\frac{K}{s(s+a)}}
  =\frac Ka\cdot\frac{1-\ee^{-aT}}{z-\ee^{-aT}},
\]
\[
  L(z)=\frac{K K_c}{a}\cdot\frac{1-\ee^{-aT}}{z-\ee^{-aT}}.
\]
Then $L(1)=\dfrac{K K_c}{a}$ and
\[
  e(\infty)=\frac{1}{1+L(1)}=\frac{1}{1+K K_c/a}=\frac{a}{a+K K_c}.
\]
\end{example}

\begin{boardnote}
The last line of this example on the board reads $e(\infty)=\dfrac{1}{1+L(1)}=\dfrac{K_cK}{a}$, followed by
``$K_c\downarrow\ e_{ss}\downarrow$''. The quantity $K_cK/a$ is $L(1)$, not the error. The correct result is the
one above, and the error \emph{decreases} when $K_c$ \emph{increases} (as already stated in the previous
example).
\end{boardnote}

\addedsection[Added problems: steady-state error]
\begin{problem}[Error constants of a type-1 loop]
A unity-feedback loop with $T=0.1\,\mathrm{s}$ has $L(z)=\dfrac{0.2\,(z+0.8)}{(z-1)(z-0.6)}$. Find the type,
$K_p$, $K_v$ and the steady-state error to a unit step and to a unit ramp.
\end{problem}
\begin{solution}
There is one pole at $z=1$, so $n=1$ and $K_p=L(1)=\infty$. The step error is $0$. The velocity constant is
\[
  K_v=\frac1T(z-1)L(z)\Big|_{z=1}=\frac{1}{0.1}\cdot\frac{0.2\cdot1.8}{1-0.6}=9,\qquad e_{ramp}(\infty)=\frac{1}{K_v}=0.111.
\]
The closed-loop poles are $0.7\pm0.52\,\jj$ ($|z|=0.87$), so the loop is stable, and a simulation of the
closed loop to the sampled ramp $r(k)=kT$ settles to the error $0.111$.
\end{solution}

\begin{problem}[Example~\ref{ex:ess-first} with numbers]
Take $K=2$, $a=1$, $K_c=4$ and $T=0.2\,\mathrm{s}$. Find $e(\infty)$ for a unit step and check stability.
\end{problem}
\begin{solution}
$L(1)=KK_c/a=8$, so $e(\infty)=1/(1+8)=0.111$. The closed-loop pole is
$z=\ee^{-aT}-\dfrac{KK_c}{a}\bigl(1-\ee^{-aT}\bigr)=0.8187-8\cdot0.1813=-0.631$, inside the unit circle,
so the loop is stable and the final value theorem applies. A time simulation gives $y(\infty)=0.8889$, i.e.\ an
error of $0.111$.
\end{solution}

%------------------------------------------------------------------------------
\subsection{Stability of digital control systems}\scanref{pp.\,46--48}
Stability is a basic requirement for analog and digital control systems. A digital control system is based on
samples and is updated every sampling period, and there is a possibility that the system may become unstable
between updates. This obviously makes stability analysis different in the digital case.

\subsubsection*{Definitions of stability}
\begin{definition}[Asymptotic stability]
A system is said to be asymptotically stable if its response to any initial conditions decays to zero
asymptotically in steady state, that is, the response due to the initial conditions satisfies
\[
  \lim_{k\to\infty}y(k)=0.
\]
\end{definition}
\begin{definition}[Marginal stability]
If the response remains bounded but does not decay to zero, the system is said to be marginally stable.
\end{definition}
\begin{definition}[BIBO stability]
A system is said to be BIBO (bounded-input bounded-output) stable if its response to any bounded input
remains bounded:
\[
  |u(k)|<b_u,\ \ 0<b_u<\infty\quad(\text{input}),
\]
\[
  \Longrightarrow\ |y(k)|<b_y,\ \ 0<b_y<\infty\quad(\text{output}),\qquad k=0,1,2,\dots
\]
\end{definition}
\begin{figure}[ht]
\centering
\figBiboSeq
\caption{A bounded output sequence: $|y(k)|<b_y$ for all $k$.}
\end{figure}

\subsubsection*{Stable $z$-domain pole locations}
Consider the sampled exponential and its $z$ transform,
\[
  p^{k}=\ee^{-akT},\quad k=0,1,2,\dots\ \longrightarrow\ \frac{z}{z-p},\qquad p\ \text{real or complex.}
\]
Then the time sequence for large $k$ is given by
\[
  |p|^{k}\ \to\ \begin{cases}0,&|p|<1,\\1,&|p|=1,\\\infty,&|p|>1.\end{cases}
\]
Any time sequence can be described by
\[
  f(k)=\sum_{i=1}^{n}A_i\,p_i^{k},\quad k=0,1,2,\dots\ \longrightarrow\ F(z)=\sum_{i=1}^{n}A_i\,\frac{z}{z-p_i},
\]
where $A_i$ are the partial-fraction coefficients and $p_i$ the $z$-domain poles. Hence we conclude that the
sequence is bounded if its poles lie on or inside the unit circle, and decays exponentially if its poles
lie \emph{inside} the unit circle (it grows exponentially if they lie outside).
Note that the case of repeated poles on the unit circle corresponds to an unbounded time sequence.

\begin{boardnote}
The board sentence ends ``\dots and decay exponentially if its poles lie outside the unit circle''; ``outside''
must be ``inside'' (the sequence \emph{grows} for poles outside). In the ramp example below the board
labels the sequence ``unbounded (expo.)'', but $y(k)=kT$ grows linearly, not exponentially.
\end{boardnote}

\begin{figure}[ht]
\centering
\begin{minipage}{0.4\linewidth}\centering\resizebox{\linewidth}{!}{\figPoleCircle}\end{minipage}\hfill
\begin{minipage}{0.55\linewidth}\centering
For example, the ramp function $\dfrac{1}{s^{2}}$ (two poles at the origin of the $s$ plane)
$\ \to\ \dfrac{Tz}{(z-1)^{2}}$ (two poles at $z=1$) $\ \to\ y(k)=kT$:\\[2mm]
\figRampSeq
\end{minipage}
\caption{Poles on or inside the unit circle give bounded sequences, except repeated poles on the circle.}
\end{figure}

%------------------------------------------------------------------------------
\subsection{Asymptotic stability and pole--zero cancellation}\scanref{p.\,49}
\begin{theorem}[Asymptotic stability]
In the absence of pole--zero cancellation, an LTI digital system is asymptotically stable if its poles are
inside the unit circle, and marginally stable if its poles are on the unit circle with no repeated poles on
the unit circle.
\end{theorem}

\paragraph{Pole--zero cancellation.}
\begin{itemize}
  \item Invisible but stable poles would not lead us to a wrong conclusion.
  \item Invisible but unstable poles would.
\end{itemize}

\begin{example}
Determine the asymptotic stability of the following systems.
\begin{center}
\renewcommand{\arraystretch}{1.7}
\begin{tabular}{@{}l c c l@{}}
 & $H(z)$ & poles & conclusion\\\hline
a) & $\dfrac{4(z-2)}{(z-2)(z-0.1)}$ & $2,\ 0.1$ & unstable\\
b) & $\dfrac{4(z-0.2)}{(z-0.2)(z-0.1)}$ & $0.2,\ 0.1$ & asymptotically stable\\
c) & $\dfrac{5(z-0.1)}{(z-0.2)(z-0.1)}$ & $0.1,\ 0.2$ & asymptotically stable\\
d) & $\dfrac{8(z-0.2)}{(z-0.1)(z-1)}$ & $0.1,\ 1$ & marginally stable\\
e) & $\dfrac{8(z-0.2)}{(z-0.1)(z-1)^{2}}$ & $0.1,\ 1,\ 1$ (repeated) & unstable
\end{tabular}
\end{center}
In (a) the cancelled pole $z=2$ is unstable although invisible in the reduced transfer function; in (b) and
(c) the cancelled pole is stable.
\end{example}

\begin{boardnote}
In part e) the board lists the poles as $0.2,\ 1,\ 1$. The value $0.2$ is the zero of $H(z)$; the third pole
is $0.1$, as in the table.
\end{boardnote}

%------------------------------------------------------------------------------
\subsection{BIBO stability}\scanref{pp.\,50--51}
\begin{theorem}
A discrete-time linear system is BIBO stable if and only if (iff) its impulse response sequence is
absolutely summable, that is,
\[
  \sum_{i=0}^{\infty}|h(i)|<\infty.
\]
\end{theorem}
A discrete-time linear system is BIBO stable if and only if the poles of its transfer function lie inside
the unit circle.

\begin{figure}[ht]
\centering
\resizebox{0.4\linewidth}{!}{\figBiboRegions}
\caption{BIBO stable (inside) and BIBO unstable (outside) pole regions.}
\end{figure}

There is no expression for the pole--zero cancellation here: BIBO stability is judged on the transfer
function as it stands, i.e.\ after any cancellation (a cancelled mode cannot be seen from the input and
the output).

\begin{example}
Determine the BIBO stability of the systems of the previous example.
\begin{center}
\renewcommand{\arraystretch}{1.7}
\begin{tabular}{@{}l c c l@{}}
 & $H(z)$ & poles & conclusion\\\hline
a) & $\dfrac{4(z-2)}{(z-2)(z-0.1)}\to\dfrac{4}{z-0.1}$ & $0.1$ & BIBO stable\\
b) & $\dfrac{4(z-0.2)}{(z-0.2)(z-0.1)}$ & $0.2,\ 0.1$ & BIBO stable\\
c) & $\dfrac{5(z-0.1)}{(z-0.2)(z-0.1)}$ & $0.1,\ 0.2$ & BIBO stable\\
d) & $\dfrac{8(z-0.2)}{(z-0.1)(z-1)}$ & $z=1,\ 0.1$ & BIBO unstable\\
e) & $\dfrac{8(z-0.2)}{(z-0.1)(z-1)^{2}}$ & $z_1=1,\ z_2=1,\ z_3=0.1$ & BIBO unstable
\end{tabular}
\end{center}
\end{example}

\paragraph{Summary.}
\begin{center}
\renewcommand{\arraystretch}{1.3}
\begin{tabular}{@{}ll@{}}
\multicolumn{2}{@{}l}{Without unstable pole--zero cancellation:}\\
asymptotically stable & $\equiv$ BIBO stable $\equiv$ stable\\
marginally stable & $\equiv$ BIBO unstable\\
unstable & $\equiv$ BIBO unstable
\end{tabular}
\end{center}

\begin{example}[Stability from a difference equation]
Determine the stability of the following system, where all initial conditions are zero:
\[
  y(k+2)+0.8\,y(k+1)+0.07\,y(k)=2u(k+1)+0.2\,u(k),\qquad k=0,1,2,\dots
\]
Taking the $z$ transform,
\[
  z^{2}Y(z)+0.8z\,Y(z)+0.07\,Y(z)=2z\,U(z)+0.2\,U(z)\ \Rightarrow\ Y(z)\bigl[z^{2}+0.8z+0.07\bigr]=U(z)\bigl[2z+0.2\bigr],
\]
\[
  \frac{Y(z)}{U(z)}=\frac{2(z+0.1)}{(z+0.7)(z+0.1)}=\frac{2(z+0.1)}{(z+0.7)(z+0.1)}\ \to\ z=-0.7
\]
(the pole at $-0.1$ is cancelled by the zero; it is stable anyway). The pole $z=-0.7$ is inside the unit circle:
the system is asymptotically stable, BIBO stable, i.e.\ stable.
\end{example}

\addedsection[Added problems: classifying stability]
\begin{problem}
Classify each system as asymptotically stable, marginally stable or unstable, and say whether it is BIBO stable:
(a) $H(z)=\dfrac{z-1.5}{(z-1.5)(z+0.3)}$; (b) $H(z)=\dfrac{1}{z^{2}+1}$; (c) $H(z)=\dfrac{1}{(z^{2}+1)^{2}}$.
\end{problem}
\begin{solution}
(a) After cancellation the pole is $-0.3$, so $H$ is BIBO stable; but the cancelled pole $z=1.5$ is unstable
and invisible, so the internal dynamics are unstable (the system is \emph{not} asymptotically stable). This
is the case of an \emph{unstable} pole--zero cancellation.\par
(b) The poles $\pm\jj$ are simple and on the unit circle: marginally stable, BIBO unstable (the bounded input
$\sin(\pi k/2)$ excites the resonance and the output grows).\par
(c) Repeated poles $\pm\jj$ on the unit circle: unstable (the response grows like $k$) and BIBO unstable.
\end{solution}

%------------------------------------------------------------------------------
\subsection{Stability determination: MATLAB and the Routh--Hurwitz criterion}\scanref{pp.\,52--56}
There are different methods for determining stability: MATLAB, the Routh--Hurwitz criterion and the Jury test.

\paragraph{\textcircled{1} MATLAB.}
\begin{itemize}
  \item \texttt{roots}: returns the roots of a polynomial.
  \item \texttt{pole}: returns the poles of a SISO or MIMO system.
  \item \texttt{pzmap}: returns the pole--zero map of a dynamic system.
\end{itemize}

\paragraph{\textcircled{2} Routh--Hurwitz criterion.}
\begin{itemize}
  \item It determines conditions for the stability of LTI systems.
  \item It cannot be used to investigate the stability of a discrete-time system directly.
  \item The bilinear transformation is needed to use the Routh--Hurwitz criterion for the stability analysis of
        discrete-time systems:
        \begin{equation}\label{eq:bilinear}
          z=\frac{1+w}{1-w}\quad\longleftrightarrow\quad w=\frac{z-1}{z+1}.
        \end{equation}
        It transforms the inside of the unit circle to the $w$-domain left half plane (LHP).
\end{itemize}

\subsubsection*{Verification of the bilinear transformation}
Since $w=\dfrac{z-1}{z+1}$, the angle of $w$ is $\angle w=\angle(z-1)-\angle(z+1)$.
\begin{figure}[ht]
\centering
\resizebox{0.98\linewidth}{!}{\figBilinear}
\caption{Geometric check of $w=(z-1)/(z+1)$ for a point inside ($z_1$), on ($z_2$) and outside ($z_3$) the unit
circle.}\label{fig:bilinear}
\end{figure}
\begin{align*}
  \angle w_1&=\angle(z_1-1)-\angle(z_1+1)\ \to\ |\angle w_1|>90^{\circ}\quad(w_1\ \text{in the LHP}),\\
  \angle w_2&=\angle(z_2-1)-\angle(z_2+1)\ \to\ \angle w_2=\pm90^{\circ}\quad(w_2\ \text{on the imaginary axis}),\\
  \angle w_3&=\angle(z_3-1)-\angle(z_3+1)\ \to\ |\angle w_3|<90^{\circ}\quad(w_3\ \text{in the RHP}).
\end{align*}
For the point $z_2$ on the circle the two segments to $\pm1$ are perpendicular (a diameter subtends a right
angle), which gives $\angle w_2=\pm90^{\circ}$.

\begin{boardnote}
For $z_2$ the board writes $\angle w_2>\pm90^{\circ}$. For a point on the circle the angle is \emph{equal} to
$\pm90^{\circ}$, as typeset above. (For $z_1$ and $z_3$ the board's $\angle w_1>90^{\circ}$ and
$\angle w_3<90^{\circ}$ are kept; they refer to the magnitude of the angle.)
\end{boardnote}

For the general $z$-polynomial
\[
  F(z)=a_nz^{n}+a_{n-1}z^{n-1}+\cdots+a_0\ \xrightarrow{\ z=\frac{1+w}{1-w}\ }\
  a_n\Bigl(\frac{1+w}{1-w}\Bigr)^{n}+a_{n-1}\Bigl(\frac{1+w}{1-w}\Bigr)^{n-1}+\cdots+a_0.
\]
\begin{itemize}
  \item For low-order polynomials it easily gives stability conditions.
  \item For high-order polynomials a symbolic manipulation package can be used, such as the MATLAB Symbolic
        Toolbox, Mathematica, Maxima (open source), etc.
\end{itemize}

\begin{example}[Second-order polynomial]
Test the stability of the polynomial $F(z)=z^{2}-0.25=0$ by using the Routh--Hurwitz criterion. With the
bilinear transformation $z=\dfrac{1+w}{1-w}$,
\[
  \Bigl(\frac{1+w}{1-w}\Bigr)^{2}-0.25=0\ \to\ 0.75\,w^{2}+2.5\,w+0.75=0.
\]
Routh table:
\[
\begin{array}{c|cc}
w^{2}&0.75&0.75\\
w^{1}&2.5&0\\
w^{0}&0.75&
\end{array}
\]
There is no sign change in the first column, so the system is stable (the roots of $F$ are $\pm0.5$).
\end{example}

\begin{example}[Third-order polynomial]
Test the stability of the polynomial $F(z)=z^{3}+1.25z^{2}-1.375z+0.25=0$ by using the Routh--Hurwitz criterion
(see the board note below for this polynomial). The bilinear transformation $z=\dfrac{1+w}{1-w}$ gives
\[
  -1.875\,w^{3}+3.875\,w^{2}+4.875\,w+1.125=0.
\]
Routh table:
\[
\begin{array}{c|cc}
w^{3}&-1.875&4.875\\
w^{2}&3.875&1.125\\
w^{1}&5.419&0\\
w^{0}&1.125&
\end{array}
\]
There is one sign change in the first column, so the system is unstable. It means that one root is in the
$w$-domain RHP. This implies that there will be one root of the characteristic equation outside the unit
circle in the $z$ plane (indeed $F(z)=(z+2)(z-\tfrac12)(z-\tfrac14)$ has the root $z=-2$).
\end{example}

\begin{boardnote}
The polynomial printed on the board is $F(z)=z^{3}-1.2z^{2}-1.375z-0.25$, but the $w$-polynomial and the Routh
table written below it belong to $F(z)=z^{3}+1.25z^{2}-1.375z+0.25=(z+2)(z-\tfrac12)(z-\tfrac14)$. The
example above is typeset with the polynomial that matches the board's calculation. For the printed
polynomial the bilinear transformation gives $1.075\,w^{3}+4.825\,w^{2}+3.925\,w-1.825=0$, whose Routh column
$1.075,\ 4.825,\ 4.332,\ -1.825$ also has one sign change (one root outside the circle, at $z=1.965$), so the
conclusion on the board is the same.
\end{boardnote}

\subsubsection*{Generalization for low-order polynomials}
\paragraph{First-order polynomial.} Consider $a_1z+a_0$ with $a_1>0$. The root is $z=-\dfrac{a_0}{a_1}$ and the
stability condition is $\left|\dfrac{a_0}{a_1}\right|<1$.

\paragraph{Second-order polynomial.} Consider $a_2z^{2}+a_1z+a_0$ with $a_2>0$. The bilinear transformation
$z=\dfrac{1+w}{1-w}$ gives
\[
  a_2\Bigl(\frac{1+w}{1-w}\Bigr)^{2}+a_1\Bigl(\frac{1+w}{1-w}\Bigr)+a_0
  =(a_2-a_1+a_0)\,w^{2}+2(a_2-a_0)\,w+(a_2+a_1+a_0).
\]
By the Routh--Hurwitz criterion, it can be shown that the poles of the second-order $w$-polynomial remain in
the LHP if and only if its coefficients are all positive. The stability condition is given by
\begin{enumerate}
  \item $a_2-a_1+a_0>0$,
  \item $a_2-a_0>0$,
  \item $a_2+a_1+a_0>0$.
\end{enumerate}

\begin{example}[Second-order loop]
Use the Routh--Hurwitz criterion to investigate the stability of the system
$G(z)=\dfrac{5(z-1)}{(z-0.1)(z-0.8)}=\dfrac{5(z-1)}{z^{2}-0.9z+0.08}$. Here $a_2z^{2}+a_1z+a_0=z^{2}-0.9z+0.08$ with $a_2=1>0$:
\begin{enumerate}
  \item $a_2-a_1+a_0=1+0.9+0.08=1.98>0$,
  \item $a_2-a_0=1-0.08=0.92>0$,
  \item $a_2+a_1+a_0=1-0.9+0.08=0.18>0$.
\end{enumerate}
All conditions are met, so the system is stable.
\end{example}

\begin{example}[Stable range of a proportional gain]
Use the Routh--Hurwitz criterion to find the stable range of the proportional controller $K$ for the
closed-loop unity-feedback system with loop gain
\[
  L(z)=K\,G(z)=\frac{K(z-1)}{(z-0.1)(z-0.8)}.
\]
The closed-loop transfer function is $G_{cl}=\dfrac{L(z)}{1+L(z)}$ and the characteristic equation is $1+L(z)=0$:
\[
  1+\frac{K(z-1)}{(z-0.1)(z-0.8)}=0\ \to\ (z-0.1)(z-0.8)+K(z-1)=0\ \to\ z^{2}+(K-0.9)\,z+0.08-K=0.
\]
We have $a_2=1>0$ \checkmark. For a second-order system the Routh--Hurwitz stability conditions give
\begin{align*}
  1.\ \ a_2-a_1+a_0>0&:\ 1-(K-0.9)+0.08-K>0\ \to\ K<0.99,\\
  2.\ \ a_2-a_0>0&:\ 1-(0.08-K)>0\ \to\ K>-0.92,\\
  3.\ \ a_2+a_1+a_0>0&:\ 1+(K-0.9)+0.08-K>0\ \to\ 0.18>0\ \checkmark
\end{align*}
Hence the closed loop is stable for $-0.92<K<0.99$.
\end{example}

\addedsection[Added problems: Routh--Hurwitz with the bilinear transformation]
\begin{problem}
Apply the second-order conditions to $F(z)=z^{2}-1.2z+0.5$ and verify with the roots.
\end{problem}
\begin{solution}
$a_2-a_1+a_0=1+1.2+0.5=2.7>0$; $a_2-a_0=0.5>0$; $a_2+a_1+a_0=1-1.2+0.5=0.3>0$: stable. The roots are
$0.6\pm0.374\,\jj$ with $|z|=\sqrt{0.5}=0.707<1$.
\end{solution}

\begin{problem}[Stable range of a gain]
For $L(z)=\dfrac{K(z+0.5)}{(z-1)(z-0.5)}$ find the range of $K$ for which the unity-feedback loop is stable.
\end{problem}
\begin{solution}
$1+L=0\Rightarrow(z-1)(z-0.5)+K(z+0.5)=z^{2}+(K-1.5)z+(0.5+0.5K)=0$. The conditions give
$a_2-a_1+a_0=3-0.5K>0\Rightarrow K<6$; $a_2-a_0=0.5-0.5K>0\Rightarrow K<1$;
$a_2+a_1+a_0=1.5K>0\Rightarrow K>0$. Hence $0<K<1$. At $K=1$ the polynomial is $z^{2}-0.5z+1$, whose roots
have $|z|=1$ (the stability limit); for $K=1.01$ they have $|z|=1.0025$.
\end{solution}

%------------------------------------------------------------------------------
\subsection{Jury test}\scanref{pp.\,57--60}
It is possible to investigate the stability of a $z$-domain polynomial directly using the Jury test. For the
polynomial
\[
  F(z)=a_nz^{n}+a_{n-1}z^{n-1}+\cdots+a_1z+a_0,\qquad a_n>0,
\]
the Jury table is

\begin{center}
\tabJury
\end{center}

The roots of the polynomial are inside the unit circle if and only if
\begin{enumerate}
  \item $F(1)>0$,
  \item $(-1)^{n}F(-1)>0$,
  \item $|a_0|<a_n$,
  \item $|b_0|>|b_{n-1}|$,
  \item $|c_0|>|c_{n-2}|$,
  \item[\vdots]
  \item[$(n{+}1)$.] $|r_0|>|r_2|$,
\end{enumerate}
where the terms in the $(n+1)$ conditions are calculated from the Jury table. The entries of the table are
calculated as follows:
\[
  b_k=\begin{vmatrix}a_0&a_{n-k}\\a_n&a_k\end{vmatrix},\ k=0,1,\dots,n-1;\qquad
  c_k=\begin{vmatrix}b_0&b_{n-1-k}\\b_{n-1}&b_k\end{vmatrix},\ k=0,1,\dots,n-2;
\]
\[
  r_0=\begin{vmatrix}s_0&s_3\\s_3&s_0\end{vmatrix},\qquad
  r_1=\begin{vmatrix}s_0&s_2\\s_3&s_1\end{vmatrix},\qquad
  r_2=\begin{vmatrix}s_0&s_1\\s_3&s_2\end{vmatrix}.
\]

\paragraph{Example: the entries for $n=5$.}
\[
  b_0=\begin{vmatrix}a_0&a_5\\a_5&a_0\end{vmatrix},\
  b_1=\begin{vmatrix}a_0&a_4\\a_5&a_1\end{vmatrix},\
  b_2=\begin{vmatrix}a_0&a_3\\a_5&a_2\end{vmatrix},\
  b_3=\begin{vmatrix}a_0&a_2\\a_5&a_3\end{vmatrix},\
  b_4=\begin{vmatrix}a_0&a_1\\a_5&a_4\end{vmatrix},
\]
\[
  c_0=\begin{vmatrix}b_0&b_4\\b_4&b_0\end{vmatrix},\
  c_1=\begin{vmatrix}b_0&b_3\\b_4&b_1\end{vmatrix},\
  c_2=\begin{vmatrix}b_0&b_2\\b_4&b_2\end{vmatrix},\
  c_3=\begin{vmatrix}b_0&b_1\\b_4&b_3\end{vmatrix},
\]
\[
  d_0=\begin{vmatrix}c_0&c_3\\c_3&c_0\end{vmatrix},\
  d_1=\begin{vmatrix}c_0&c_2\\c_3&c_1\end{vmatrix},\
  d_2=\begin{vmatrix}c_0&c_1\\c_3&c_2\end{vmatrix}.
\]

\begin{example}[Fifth-order polynomial]
Test the stability of the polynomial
\[
  F(z)=z^{5}+2.6z^{4}-0.56z^{3}-2.05z^{2}+0.0775z+0.35=0,\qquad n=5\ \to\ 2n-3=7\ \text{rows.}
\]
\begin{center}
\small
\renewcommand{\arraystretch}{1.25}
\begin{tabular}{c|cccccc}
Row & $z^{0}$ & $z^{1}$ & $z^{2}$ & $z^{3}$ & $z^{4}$ & $z^{5}$\\\hline
1 & $a_0=0.35$ & $a_1=0.0775$ & $a_2=-2.05$ & $a_3=-0.56$ & $a_4=2.6$ & $a_5=1$\\
2 & $a_5=1$ & $a_4=2.6$ & $a_3=-0.56$ & $a_2=-2.05$ & $a_1=0.0775$ & $a_0=0.35$\\
3 & $b_0=-0.8775$ & $b_1=-2.5729$ & $b_2=-0.1575$ & $b_3=1.854$ & $b_4=0.8325$ & \\
4 & $b_4=0.8325$ & $b_3=1.854$ & $b_2=-0.1575$ & $b_1=-2.5729$ & $b_0=-0.8775$ & \\
5 & $c_0=0.0770$ & $c_1=0.7143$ & $c_2=0.2693$ & $c_3=0.5151$ & & \\
6 & $c_3=0.5151$ & $c_2=0.2693$ & $c_1=0.7143$ & $c_0=0.0770$ & & \\
7 & $d_0=-0.2593$ & $d_1=-0.0837$ & $d_2=-0.3472$ & & &
\end{tabular}
\end{center}
For instance $b_0=\begin{vmatrix}a_0&a_5\\a_5&a_0\end{vmatrix}=0.35^{2}-1=-0.8775$. Stability conditions:
\begin{enumerate}
  \item $F(1)=1.4175>0$ \checkmark
  \item $(-1)^{5}F(-1)=(-1)\cdot0.3825<0$ \quad\textbf{\textcolor{boardred}{violated}}
  \item $|a_0|<a_5$: $0.35<1$ \checkmark
  \item $|b_0|>|b_4|$: $0.8775>0.8325$ \checkmark
  \item $|c_0|>|c_3|$: $0.0770>0.5151$ \quad\textbf{\textcolor{boardred}{violated}}
  \item $|d_0|>|d_2|$: $0.2593<0.3472$ \quad\textbf{\textcolor{boardred}{violated}}
\end{enumerate}
A single violation of one of the conditions is enough for the system \emph{not} to be stable. The
polynomial is therefore unstable (its roots are $-2.5,\,-0.8,\,0.7,\,0.5,\,-0.5$; $|z|=2.5$ lies outside the
unit circle).
\end{example}

\begin{boardnote}
At the end of this example the board says ``(2 or 6)'' for the violated conditions and does not mark condition
5, but in the board's own numbers condition 5 ($0.0770>0.5151$) is also false (it is written without a
mark). Conditions 2, 5 and 6 are violated. Since the polynomial is unstable, the conclusion is unchanged.
\end{boardnote}

\begin{example}[Stable range of the gain, vehicle position]
Find the stable range of the proportional controller gain $K$ for the vehicle position control system with
the analog plant
\[
  G(s)=\frac{10}{s(s+10)}\quad\text{and DAC and ADC, if } T=0.05\,\mathrm{s}.
\]
\[
  \GZAS(z)=(1-z^{-1})\ZT{\frac{G(s)}{s}}=(1-z^{-1})\,\mathcal{Z}\!\left[0.1\Bigl(\frac{10}{s^{2}}-\frac1s+\frac{1}{s+10}\Bigr)\right]
\]
\[
  \Rightarrow\ \GZAS(z)=\frac{1.0653\times10^{-2}\,(z+0.8467)}{(z-1)(z-0.6065)}.
\]
The closed-loop transfer function is $G_{cl}(z)=\dfrac{K\GZAS(z)}{1+K\GZAS(z)}$ with characteristic equation
$1+K\GZAS(z)=0$:
\[
  (z-1)(z-0.6065)+1.0653\times10^{-2}K(z+0.8467)=0
\]
\[
  z^{2}+\bigl(1.0653\times10^{-2}K-1.6065\bigr)z+0.6065+9.02\times10^{-3}K=0.
\]
Jury conditions ($n=2$):
\begin{enumerate}
  \item $F(1)>0$: $1+(1.0653\times10^{-2}K-1.6065)+0.6065+9.02\times10^{-3}K>0\ \Rightarrow\ K>0$.
  \item $(-1)^{2}F(-1)>0$: $1-(1.0653\times10^{-2}K-1.6065)+0.6065+9.02\times10^{-3}K>0\ \Rightarrow\ K<1967.582$.
  \item $|a_0|<a_2$: $|0.6065+9.02\times10^{-3}K|<1\ \Rightarrow\ -178.104<K<43.6199$.
\end{enumerate}
Stability range for $K$: $0<K<43.6199$.
\end{example}

\begin{boardnote}
In condition 3 the board writes $|0.6065+0.0902K|<1$; the coefficient is $9.02\times10^{-3}=0.00902$, as in the
characteristic equation above. The bounds $-178.104<K<43.6199$ on the board are the ones obtained with $0.00902$.
\end{boardnote}

\addedsection[Added problem: Jury test, third order]
\begin{problem}
Use the Jury test on $F(z)=z^{3}-0.9z^{2}-0.22z+0.24$ and on $F(z)=z^{3}-1.2z^{2}-0.19z+0.33$.
\end{problem}
\begin{solution}
For $n=3$ the table has three rows and the conditions are $F(1)>0$, $(-1)^{3}F(-1)>0$, $|a_0|<a_3$ and
$|b_0|>|b_2|$, with $b_0=a_0^{2}-a_3^{2}$ and $b_2=a_0a_2-a_3a_1$.\par
\emph{First polynomial} ($a_0=0.24$, $a_1=-0.22$, $a_2=-0.9$, $a_3=1$): $F(1)=0.12>0$;
$-F(-1)=1.44>0$; $0.24<1$; $b_0=-0.9424$, $b_2=0.004$, so $0.9424>0.004$. All hold, the polynomial is stable
(its roots are $0.8,\,-0.5,\,0.6$).\par
\emph{Second polynomial} ($a_0=0.33$, $a_1=-0.19$, $a_2=-1.2$, $a_3=1$): $F(1)=-0.06<0$, so it is not stable
(one root is $z=1.1$; the others are $-0.5$ and $0.6$).
\end{solution}
